%========================================================
% Reviewer-guided reformulation for A10-PTF (Phase 1: RHPS)
%========================================================

\subsection{Reviewer-guided reformulation of the RHPS model}
\label{subsec:rhps_reformulation}

To obtain a mathematically closed reduced-order model suitable for optimization,
we explicitly separate prognostic states from diagnostic observables.
The reduced RHPS state is defined by
\begin{equation}
x(t)=
\begin{bmatrix}
a(t)\\
z_f(t)\\
z_N(t)\\
\eta(t)
\end{bmatrix},
\label{eq:rhps_state}
\end{equation}
where
$a(t)\in\mathbb{R}^{2N_a}$ is the retained acoustic modal state,
$z_f(t)\in\mathbb{R}^{N_f}$ is an internal flame-memory state,
$z_N(t)\in\mathbb{R}^{N_N}$ is a causal nozzle-memory state,
and $\eta(t)\in\mathbb{R}^{N_\eta}$ is a slow thermochemical state
(e.g.\ mean mixture, mean temperature, and mean mass-flow coordinates).

The reduced dynamics are written as
\begin{align}
\dot a
&=
A_a(\theta,\eta)\,a
+
B_f(\theta,\eta)\,y_f
+
B_N(\theta,\eta)\,y_N
+
B_u(\theta,\eta)\,u
+
\mathcal{N}_a(a,z_f,z_N,\eta,u;\theta),
\label{eq:rhps_a}
\\
\dot z_f
&=
A_f(\eta)\,z_f
+
B_{fa}(\eta)\,a
+
B_{fu}(\eta)\,u
+
\mathcal{N}_f(a,z_f,\eta,u;\theta),
\label{eq:rhps_zf}
\\
\dot z_N
&=
A_N(\eta)\,z_N
+
B_{Np}(\eta)\,p_c'(a)
+
\mathcal{N}_N(z_N,p_c'(a);\eta),
\label{eq:rhps_zn}
\\
\dot\eta
&=
f_\eta(\eta,\bar T,\bar\phi,\bar{\dot m};u,\theta),
\label{eq:rhps_eta}
\end{align}
with outputs
\begin{equation}
y_f = C_f(\eta)\,z_f + D_f(\eta)\,a,
\qquad
y_N = C_N(\eta)\,z_N + D_N(\eta)\,p_c'(a).
\label{eq:rhps_outputs}
\end{equation}
Here, $\theta$ denotes the uncertain operating and forcing parameters.

\paragraph{Causal nozzle-memory realization.}
Instead of a formal nonlinear convolution written directly at the boundary,
we realize the nozzle response through the stable finite-dimensional state $z_N$.
In the linearized limit,
\begin{equation}
K_N(t;\eta)=C_N(\eta)\,e^{A_N(\eta)t}\,B_{Np}(\eta),
\qquad
\Re \lambda_i(A_N)<0
\ \ \forall i,
\label{eq:nozzle_kernel}
\end{equation}
so that the associated kernel is causal and exponentially decaying.
This gives a realizable boundary memory model while preserving a state-space form
for optimization and adjoint differentiation.

\subsection{TDHR as a regularized diagnostic observable}
\label{subsec:tdhr_regularized}

The proposed topological discontinuity of heat release (TDHR) is not introduced as a
literal discontinuity or as a prognostic state.
Instead, it is treated as a \emph{regularized diagnostic field} defined on the flame sheet
or on a thin reaction zone $\Gamma_f(t)$.
Let $\dot q_\sigma$ denote a filtered heat-release field and let
\begin{equation}
S_\kappa(s)=\frac{1}{2}\left(1+\tanh(\kappa s)\right)
\label{eq:smooth_switch}
\end{equation}
be a smooth switching function.
We define the TDHR indicator
\begin{equation}
\chi_{\mathrm{TDHR}}(x,t)
=
S_\kappa\!\left(\|\nabla_{\Gamma}\dot q_\sigma\|-\xi_q\right)
\,
S_\kappa\!\left(\xi_H-\det H_{\Gamma}(\dot q_\sigma)\right),
\label{eq:tdhr_indicator}
\end{equation}
where $\nabla_{\Gamma}$ and $H_{\Gamma}$ are the surface gradient and surface Hessian
restricted to $\Gamma_f(t)$, respectively, and $\xi_q,\xi_H>0$ are detection thresholds.

The retained modal TDHR diagnostic is then defined by
\begin{equation}
D_n^{\mathrm{TDHR}}(t)
=
\int_{\Gamma_f(t)}
p_n^\dagger(x)\,
\chi_{\mathrm{TDHR}}(x,t)\,
\dot q'(x,t)\,dA,
\qquad n=1,\dots,N_a,
\label{eq:tdhr_modal}
\end{equation}
where $p_n^\dagger$ is the adjoint pressure (or adjoint acoustic potential)
associated with the $n$th retained thermoacoustic mode.
Equation~\eqref{eq:tdhr_modal} ensures that TDHR is interpreted only as a
mode-weighted precursor diagnostic.
It is \emph{not} taken as the primary stability functional.

To monitor whether TDHR is actually informative, we define the mismatch
\begin{equation}
\varepsilon_{\mathrm{TDHR},n}
=
\left|
D_n^{\mathrm{TDHR}}
-
\mathcal{R}_n^\dagger
\right|,
\qquad
\mathcal{R}_n^\dagger
=
\frac{1}{T_p}
\int_0^{T_p}
\int_{\Omega_f}
p_n^\dagger(x)\,\dot q'(x,t)\,dV\,dt,
\label{eq:tdhr_mismatch}
\end{equation}
where $\mathcal{R}_n^\dagger$ is the period-averaged adjoint-weighted Rayleigh source term.
If $\varepsilon_{\mathrm{TDHR},n}$ remains large over the training or validation set,
TDHR must be interpreted as a weak correlate rather than a control-relevant indicator.

\subsection{Replacement of the thermodynamic line functional}
\label{subsec:replace_Fth}

The line functional
\begin{equation}
\mathcal{F}_{th}=\oint_\Gamma \left(\dot S_{gen}-\Psi_{sub}\right)\,d\gamma
\end{equation}
is not adopted as a stability functional.
Instead, entropy generation and subgrid penalties are separated and assigned distinct roles:
\begin{align}
\mathcal{S}_{gen}(u;\theta)
&=
\int_0^T\int_{\Omega}
\dot s_{gen}(x,t;u,\theta)\,dV\,dt,
\label{eq:Sgen}
\\
\mathcal{P}_{sub}(u;\theta)
&=
\int_0^T
\Psi_{sub}(t;u,\theta)\,dt.
\label{eq:Psub}
\end{align}
The first is a thermodynamic regularizer; the second is an explicit model-discrepancy penalty.
Neither is identified with the thermoacoustic growth rate itself.

The optimization objective is written as
\begin{equation}
J(u)
=
(1-\lambda)\,\mathbb{E}_{\theta}\!\left[\mathcal{J}(u;\theta)\right]
+
\lambda\,\mathrm{CVaR}_\beta\!\left(\mathcal{J}(u;\theta)\right),
\label{eq:outer_risk}
\end{equation}
with sample cost
\begin{align}
\mathcal{J}(u;\theta)
&=
\int_0^T
\Big[
w_p \|p'(t)\|_2^2
+
w_q \|\dot q'(t)\|_2^2
+
w_{sub}\,\Psi_{sub}(t)
+
w_S\,\dot{\mathcal{S}}_{gen}(t)
+
w_u \|u(t)\|_2^2
+
w_{\dot u}\|\dot u(t)\|_2^2
\Big]\,dt
\nonumber\\
&\quad
+
w_{\mathrm{fail}}
\int_0^T
\mathcal{H}_\epsilon\!\left(g_{\mathrm{fail}}(x(t),\theta)\right)\,dt,
\label{eq:sample_cost}
\end{align}
where $\mathcal{H}_\epsilon$ is a smooth approximation of the Heaviside function
and $g_{\mathrm{fail}}$ encodes mission-level failure metrics
(e.g.\ pressure exceedance, wall-heat exceedance, or blowout proximity).

Thermoacoustic stabilization is enforced not by minimizing
$\mathcal{S}_{gen}$ directly, but by explicit stability constraints.
For an equilibrium or weakly nonlinear operating point,
\begin{equation}
\Re\lambda_j\!\left(\mathcal{L}(x^\star,u,\theta)\right)\le -\delta_j,
\qquad
j=1,\dots,N_r,
\label{eq:stability_constraint_linear}
\end{equation}
where $\mathcal{L}$ is the linearized RHPS operator and $\delta_j>0$ is a prescribed margin.
Near a periodic orbit, this is replaced by the adjoint-weighted acoustic-energy condition
\begin{equation}
\overline{\mathcal{R}_j^\dagger}
\le
\overline{\mathcal{D}_j}
-\delta_j,
\qquad
j=1,\dots,N_r,
\label{eq:stability_constraint_periodic}
\end{equation}
where $\overline{\mathcal{D}_j}$ denotes the corresponding averaged damping contribution.

\subsection{Adjoint-consistent differentiation and numerical admissibility}
\label{subsec:adjoint_admissibility}

Gradients are computed from the \emph{fully discretized} RHPS model.
Let the one-step discretization be
\begin{equation}
x_{k+1}=\Phi_{\Delta t}(x_k,u_k,\theta).
\label{eq:discrete_map}
\end{equation}
The discrete adjoint is then
\begin{equation}
\lambda_k
=
\left(\partial_{x_k}\Phi_{\Delta t}\right)^\top \lambda_{k+1}
+
\partial_{x_k}\ell_k,
\qquad
\lambda_N=\partial_{x_N}\ell_N,
\label{eq:discrete_adjoint}
\end{equation}
with gradient
\begin{equation}
\frac{dJ}{du_k}
=
\partial_{u_k}\ell_k
+
\left(\partial_{u_k}\Phi_{\Delta t}\right)^\top \lambda_{k+1}.
\label{eq:discrete_gradient}
\end{equation}

To preserve differentiability, the following replacements are mandatory:
\begin{enumerate}
\item hard thresholds $\rightarrow$ smooth switches $S_\kappa(\cdot)$,
\item $\max(\cdot)$ $\rightarrow$ log-sum-exp or softplus approximations,
\item event counts $\rightarrow$ mollified event densities,
\item absolute values $\rightarrow$ $(x^2+\epsilon^2)^{1/2}$ regularizations where needed.
\end{enumerate}

The discrete adjoint is accepted only if it passes a Taylor remainder test of the form
\begin{equation}
\left|
J(u+\alpha \delta u)-J(u)-\alpha \nabla J(u)^\top \delta u
\right|
=
\mathcal{O}(\alpha^2)
\qquad
(\alpha\to 0).
\label{eq:taylor_test}
\end{equation}

\paragraph{Adjoint breakdown criterion.}
The standard backward adjoint is not used indiscriminately.
It is declared numerically inadmissible whenever at least one of the following holds:
\begin{align}
\max_{1\le j\le N_r}\Re\lambda_j(\mathcal{L}) &\ge -\epsilon_H,
\label{eq:hopf_margin}
\\
T_{\mathrm{opt}} &> c_\Lambda\,\Lambda_1^{-1},
\label{eq:lyapunov_limit}
\\
\mathrm{cond}\!\left(I-\partial_x\Phi_T\right) &> \kappa_{\max},
\label{eq:condition_number_limit}
\end{align}
where $\epsilon_H>0$ is a Hopf-margin tolerance,
$\Lambda_1$ is the maximal finite-time Lyapunov exponent of the reduced model,
and $c_\Lambda=\mathcal{O}(1)$ is a chosen safety constant.

If any of \eqref{eq:hopf_margin}--\eqref{eq:condition_number_limit} is triggered,
the optimization is switched from a global unsteady adjoint to one of the following:
\begin{enumerate}
\item a short-horizon receding-horizon adjoint,
\item a periodic-orbit continuation / Floquet sensitivity calculation,
\item a shadowing-type sensitivity estimator for long-time averaged objectives.
\end{enumerate}

\subsection{Claim hierarchy}
\label{subsec:claim_hierarchy_rhps}

The revised RHPS formulation is deliberately modest in scope.
It does not claim first-principles closure of the full reactive Navier--Stokes system.
It claims a control-oriented, uncertainty-aware, and adjoint-compatible
intermediate surrogate in which
\begin{enumerate}
\item local topological flame events are treated diagnostically rather than prognostically,
\item stability is enforced through modal growth or adjoint-weighted energy constraints,
\item nozzle history dependence is represented by a realizable causal memory state,
\item the optimization objective is risk-aware and explicitly penalizes subgrid-model fragility.
\end{enumerate}
Within this restricted scope, RHPS becomes mathematically better posed for optimization,
calibration, and uncertainty propagation than the original TDHR--$\mathcal{F}_{th}$ formulation.

\begin{thebibliography}{99}

\bibitem{YoshimuraUAPST2026}
K.~Yoshimura,
\newblock \emph{Uncertainty-Aware Planetary Shield Theory (UAPST):
A Closed Intermediate Surrogate for Tail-Risk-Aware Hybrid Control of Planetary Magnetic Shielding},
\newblock independent manuscript, 2026.

\bibitem{Magri2019}
L.~Magri, M.~P.~Juniper, and J.~Moeck,
\newblock Sensitivity of the Rayleigh criterion in thermoacoustics,
\newblock \emph{Journal of Fluid Mechanics}, \textbf{882}, R1, 2020.
\newblock doi:10.1017/jfm.2019.860

\bibitem{Durox2009}
D.~Durox, T.~Schuller, N.~Noiray, A.~L.~Birbaud, and S.~Candel,
\newblock Rayleigh criterion and acoustic energy balance in unconfined self-sustained oscillating flames,
\newblock \emph{Combustion and Flame}, \textbf{156}(1), 106--119, 2009.
\newblock doi:10.1016/j.combustflame.2008.07.016

\bibitem{Juniper2018}
M.~P.~Juniper,
\newblock Sensitivity Analysis of Thermoacoustic Instability with Adjoint Helmholtz Solvers,
\newblock \emph{Physical Review Fluids}, \textbf{3}(11), 110509, 2018.
\newblock doi:10.1103/PhysRevFluids.3.110509

\bibitem{Illingworth2012}
S.~J.~Illingworth, I.~C.~Waugh, and M.~P.~Juniper,
\newblock Finding thermoacoustic limit cycles for a ducted Burke-Schumann flame,
\newblock \emph{Proceedings of the Combustion Institute}, \textbf{34}(1), 911--920, 2013.
\newblock doi:10.1016/j.proci.2012.06.017

\bibitem{Waugh2013}
I.~Waugh, S.~Illingworth, and M.~Juniper,
\newblock Matrix-free continuation of limit cycles for bifurcation analysis of large thermoacoustic systems,
\newblock \emph{Journal of Computational Physics}, \textbf{240}, 225--247, 2013.
\newblock doi:10.1016/j.jcp.2012.12.034

\bibitem{Nguyen1988}
T.~V.~Nguyen,
\newblock \emph{Computer Code for the Prediction of Nozzle Admittance},
\newblock technical report N90-28630, Aerojet TechSystems Company, 1988.

\bibitem{Padmanabhan1975}
M.~S.~Padmanabhan, E.~A.~Powell, and B.~T.~Zinn,
\newblock \emph{Effect of Nozzle Nonlinearities Upon Nonlinear Stability of Liquid Propellant Rocket Motors},
\newblock NASA CR-134880, 1975.

\bibitem{Ni2019}
A.~Ni, Q.~Wang, P.~Fernandez, and C.~Talnikar,
\newblock Sensitivity analysis on chaotic dynamical systems by Finite Difference Non-Intrusive Least Squares Shadowing (FD-NILSS),
\newblock \emph{Journal of Computational Physics}, \textbf{394}, 615--631, 2019.
\newblock doi:10.1016/j.jcp.2019.06.004

\end{thebibliography}